% called by main.tex
%
\chapter{Análisis Comparativo de Resultados}
\label{ch::chapter9}

Los capítulos anteriores han presentado los experimentos de forma progresiva: primero la línea base y el ajuste supervisado serializado, después las variantes con predicción explícita
de nivel y, finalmente, la construcción de preferencias para la etapa de DPO. Este capítulo reúne esos resultados desde una perspectiva comparativa.

Como ya he introducido, el problema estudiado no consiste únicamente en producir texto parecido a
YAML, sino en generar una representación intermedia que pueda reconstruirse de forma determinista como un manifiesto Kubernetes. Por tanto, las métricas más informativas son aquellas
que atraviesan esa frontera: la parseabilidad de la salida estructurada, la parseabilidad del YAML reconstruido, la precisión de la variable \texttt{level}, la adecuación al prompt y
las señales aproximadas de validez de dominio. Las métricas de similitud textual o de igualdad exacta con la referencia se conservan como contexto, pero no constituyen el criterio
principal de éxito.

%----------------------------------------------------------------------
\section{Criterios de comparación}
\label{sec:ch9_criteria}
%----------------------------------------------------------------------

La comparación se organiza alrededor de tres preguntas. La primera es si el ajuste supervisado sobre una superficie estructurada mejora de forma clara la generación zero-shot. Esta
pregunta corresponde a la transición entre la línea base y \texttt{serialized\_sft}. La segunda es si la jerarquía debe permanecer serializada como texto o si conviene predecirla
mediante una cabeza estructural explícita. Esta es la comparación central entre \texttt{serialized\_sft} y las variantes \texttt{two\_head}. La tercera pregunta aparece después del
SFT: si un modelo ya estable en la superficie de bloques puede refinarse mediante preferencias sin perder la regularidad estructural aprendida.

La Tabla~\ref{tab:ch9_validation_summary} resume las métricas principales disponibles en validación para las familias más representativas.

\begin{table}[htbp]
\centering
\caption{Resumen comparativo en validación de los modelos principales. Los valores de DPO corresponden a la run \texttt{beta=0.10}; la variante DPO v2 se mantiene fuera de la tabla
         hasta completar su validación final.}
\label{tab:ch9_validation_summary}
\resizebox{\textwidth}{!}{%
\begin{tabular}{lccccc}
\toprule
\textbf{Métrica} & \textbf{Baseline} & \textbf{SFT} & \textbf{THM v1} & \textbf{THOP-FiLM} & \textbf{DPO 0.10} \\
\midrule
\texttt{yaml\_parse\_success\_rate}      & 30{,}77 & 98{,}57 & 48{,}57 & 28{,}57 & 97{,}14 \\
\texttt{average\_line\_text\_f1}         & 20{,}86 & 82{,}06 & 37{,}43 &  7{,}82 & 80{,}92 \\
\texttt{average\_level\_exact\_match}    & 11{,}31 & 75{,}78 & 35{,}04 &  2{,}78 & 74{,}22 \\
\texttt{average\_level\_mae}             & 0{,}7896 & 0{,}2723 & 0{,}2392 & 0{,}8430 & 0{,}2630 \\
\texttt{average\_prompt\_requirement\_f1} & 22{,}13 & 85{,}31 & 42{,}06 &  7{,}99 & 83{,}68 \\
\texttt{average\_semantic\_key\_f1}       & 27{,}78 & 95{,}52 & 45{,}92 & 10{,}13 & 94{,}06 \\
\texttt{kubernetes\_domain\_score}        & -- & 83{,}10 & 42{,}38 & 12{,}50 & 82{,}86 \\
\texttt{kubernetes\_gate\_pass\_rate}     & -- & 14{,}29 & 12{,}86 &  0{,}00 & 12{,}86 \\
\bottomrule
\end{tabular}
}
\end{table}

La tabla muestra una primera lectura bastante nítida: el salto más grande del proyecto ocurre entre la línea base y el SFT serializado. Dentro de la rama con predicción explícita de
nivel, el mejor punto global sigue siendo el THM v1. Las variantes ordinales y posicionales posteriores no deben entenderse como mejoras sucesivas de rendimiento, sino como intentos
de explicar y corregir el fallo específico de esa primera arquitectura en los niveles profundos. DPO, por su parte, conserva buena parte del comportamiento del SFT, pero no lo mejora
de forma global en su primera configuración completa.

%----------------------------------------------------------------------
\section{De la línea base al SFT serializado}
\label{sec:ch9_baseline_sft}
%----------------------------------------------------------------------

La línea base zero-shot ya permite observar una separación que recorre toda la memoria: seguir un formato textual no equivale a mantener una estructura jerárquica válida. En validación,
el baseline produce una salida estructurada parseable en la mayoría de los casos, pero solo el 30{,}77\,\% de las muestras se reconstruye como YAML válido. La diferencia entre ambos
valores indica que el modelo entiende en parte la instrucción de formato, pero no conserva de forma robusta la relación entre niveles, listas y mappings a lo largo de toda la secuencia.

El \texttt{serialized\_sft} cambia esa situación de manera sustancial. Al entrenar el modelo directamente sobre \texttt{blocks\_tsv\_v1}, la jerarquía se incorpora a la superficie
generada mediante la columna \texttt{level}. A primera vista, esta decisión podría parecer poco sofisticada: el nivel sigue siendo texto y no una variable estructural separada. Sin
embargo, los resultados del Capítulo~\ref{ch::chapter5} muestran que esta formulación es muy efectiva. La tasa de parseabilidad YAML sube hasta el 98{,}57\,\%, la coincidencia exacta
de nivel alcanza el 75{,}78\,\% y las métricas de prompt y claves semánticas mejoran de forma igualmente marcada.

La lectura relevante no es solo que el SFT mejore al baseline, algo esperable en una tarea supervisada, sino que la representación elegida encaja bien con el parser. El modelo aprende
a emitir una superficie que contiene suficiente información estructural para reconstruir el YAML casi siempre. Dicho de otro modo, el SFT serializado no resuelve la jerarquía porque
entienda formalmente el árbol YAML, sino porque aprende una convención textual que hace visible esa jerarquía para el módulo determinista posterior. Esa combinación entre superficie
aprendida y parser explica por qué \texttt{serialized\_sft} se convierte en la referencia fuerte del resto del trabajo.

También conviene conservar una limitación de esta etapa. El buen resultado estructural no implica corrección Kubernetes completa. Como se discutió en el Capítulo~\ref{ch::chapter5},
el \texttt{kubernetes\_domain\_gate\_pass\_rate} permanece en el 14{,}29\,\%, principalmente por antipatrones de seguridad y buenas prácticas que no impiden el parseo YAML. Por tanto,
el SFT serializado demuestra que la superficie estructurada es aprendible, pero no agota la evaluación de dominio.

%----------------------------------------------------------------------
\section{Jerarquía serializada frente a jerarquía explícita}
\label{sec:ch9_sft_twohead}
%----------------------------------------------------------------------

La comparación entre \texttt{serialized\_sft} y \texttt{two\_head\_sft} es la comparación conceptual más importante del trabajo. En la primera formulación, el nivel se genera como un
campo textual ordinario dentro de la secuencia. En la segunda, el texto de la línea y la profundidad jerárquica se separan: el modelo genera \texttt{content\_blocks\_v1} y una cabeza
adicional predice \texttt{level} a partir de estados latentes del transformer. Esta separación responde directamente a la hipótesis de que la jerarquía podría modelarse mejor como una
señal estructural explícita.

La comparación global en validación favorece claramente al SFT serializado. Frente al 98{,}57\,\% de YAML parseable de \texttt{serialized\_sft}, el primer \texttt{two\_head\_sft}
alcanza el 48{,}57\,\%. También cae el F1 de texto de línea, la coincidencia exacta de nivel, la adecuación al prompt y la cobertura semántica aproximada. Si se considerase únicamente
esta tabla agregada, la lectura sería sencilla: la serialización de \texttt{level} como texto funciona mejor que la cabeza explícita.

Sin embargo, esa lectura sería incompleta. El análisis condicionado a salidas parseables muestra que una parte importante de la caída global procede del cuello de botella
parser-facing. Cuando \texttt{two\_head\_sft} sí produce YAML parseable, sus métricas se acercan mucho más a las de \texttt{serialized\_sft}. En los mismos 34 ejemplos donde la rama
two-head cruza el parser, el F1 de texto queda relativamente próximo, la métrica de claves semánticas es prácticamente equivalente y el error absoluto medio de nivel favorece
ligeramente al modelo con cabeza explícita. Esta vista no invalida la comparación global, porque las salidas no parseables siguen siendo fallos reales, pero matiza qué tipo de fallo
se está observando.

El problema principal de \texttt{two\_head\_sft} no es que genere siempre contenido irreconocible, sino que la frontera entre contenido generado, nivel predicho y reconstrucción
determinista se vuelve frágil. Los errores se concentran en listas y mappings anidados, especialmente en estructuras con \texttt{containers}, \texttt{command}, \texttt{volumeMounts}
y \texttt{volumes}. En esos casos, desplazar una línea uno o dos niveles no produce solo un error local, sino que puede hacer que el parser interprete una lista como un mapping, un
campo como un elemento de secuencia o un fragmento de comando como una clave YAML mal formada.

Esta diferencia permite formular una lectura más justa de la comparación principal. Como sistema final, \texttt{serialized\_sft} es superior en el estado actual del trabajo. Como
experimento de investigación, \texttt{two\_head\_sft} revela algo que el SFT serializado ocultaba: separar explícitamente la jerarquía del texto no basta si el estado latente elegido
no contiene una señal suficientemente estable para los niveles profundos y para los cambios locales de estructura.

%----------------------------------------------------------------------
\section{Evolución de la rama two-head}
\label{sec:ch9_twohead_evolution}
%----------------------------------------------------------------------

Los experimentos posteriores de la rama two-head pueden leerse como intentos de explicar y corregir esa fragilidad, no como una sustitución directa del THM v1. El primer resultado
relevante fue que una clasificación categórica estándar de \texttt{level} ofrecía el mejor rendimiento global de la familia, pero tendía a refugiarse en los niveles más frecuentes.
Esto era coherente con la distribución del dataset, donde los niveles 0 a 4 concentran la mayor parte de las líneas, pero era insuficiente para Kubernetes: muchas estructuras con
listas internas requieren niveles 5, 6, 7 u 8.

La reformulación ordinal introdujo una hipótesis más adecuada: los niveles no son clases independientes, sino valores ordenados. Esta decisión permitió por primera vez activar niveles
profundos en algunas variantes, pero también hizo visible otro problema. La trayectoria del score ordinal no dependía solo de la estructura YAML, sino también del momento de la
generación. El modelo parecía aprender una relación fuerte entre posición temporal y profundidad, una relación parcialmente útil porque los documentos suelen comenzar por niveles
superficiales y profundizar más adelante, pero insuficiente para representar la jerarquía real.

El análisis por cuartiles del YAML hizo especialmente visible este fenómeno. En las primeras partes de la secuencia, la cabeza tendía a predecir niveles demasiado superficiales; en
las zonas intermedias y finales aparecía una rampa de profundidad más clara, pero los niveles profundos seguían comprimidos. La hipótesis de deriva posicional latente encaja con ese
patrón: el mismo nivel jerárquico no se presenta necesariamente igual en el espacio oculto al principio, en la mitad o al final de la generación.

La variante THOP-FiLM se diseñó precisamente para introducir la posición como una señal de condicionamiento más rica. Frente a la concatenación posicional tardía, FiLM modula
representaciones intermedias del cabezal ordinal. El resultado mejora la parseabilidad respecto a las variantes ordinal-posicionales anteriores y reduce algunos errores de mapping,
pero no resuelve el problema completo: en validación final, el modelo sigue mostrando una compresión fuerte de niveles profundos y una tasa de parseabilidad muy inferior tanto a la del
SFT serializado como a la del THM v1. Por tanto, THOP-FiLM no debe presentarse como una arquitectura ganadora, sino como un diagnóstico adicional: la posición importa, pero puede
convertirse en un atajo si no se separa con cuidado de la profundidad estructural.

Esta evolución deja una enseñanza metodológica para el análisis comparativo. La rama two-head no fracasa por una única causa simple. Su comportamiento combina desbalanceo de niveles,
calibración ordinal, elección del hidden state, deriva posicional y sensibilidad del parser ante errores locales de indentación. La secuencia experimental sugiere una lectura concreta:
el THM v1 era el generador explícito más estable, pero ocultaba un fallo importante en niveles profundos; al atacar ese fallo con umbrales ordinales se hizo visible una nueva dificultad,
la no estacionariedad de la señal latente a lo largo de la generación. FiLM ayuda a formular esa dificultad, aunque todavía no la resuelve dentro de la arquitectura ordinal. Esta mezcla
explica por qué añadir una cabeza estructural es una intervención conceptualmente atractiva, pero experimentalmente más exigente que serializar la jerarquía como texto.

%----------------------------------------------------------------------
\section{Efecto de DPO sobre la referencia serializada}
\label{sec:ch9_dpo}
%----------------------------------------------------------------------

DPO parte de una situación distinta. La pregunta ya no es si el modelo puede aprender una representación estructurada básica, porque \texttt{serialized\_sft} ya lo ha demostrado. La
pregunta es si un modelo que genera bloques estructurados con alta estabilidad puede ser desplazado hacia salidas más adecuadas al prompt y más coherentes con ciertos criterios de
dominio mediante preferencias automáticas.

El primer experimento completo, con \(\beta=0{,}10\), ofrece una respuesta mixta. La parseabilidad YAML baja ligeramente respecto al SFT, del 98{,}57\,\% al 97{,}14\,\%, y también
descienden de forma leve el F1 de texto de línea, la coincidencia exacta de nivel, la adecuación al prompt y las claves semánticas. Al mismo tiempo, DPO mejora algunas señales
localizadas: sube el \texttt{parsed\_equal\_rate}, mejora ligeramente el ajuste del número de documentos y reduce ciertos errores de dominio, como \texttt{volume\_mount\_without\_volume}.
La descomposición por niveles Kubernetes muestra además una mejora en niveles intermedios de validez, aunque el gate final no mejora.

Esta combinación es importante porque muestra el riesgo de una fase de alineamiento automática. DPO no actúa como una mejora monotónica sobre el SFT, sino como una reponderación del
comportamiento del modelo. Puede empujar algunas propiedades de dominio en la dirección deseada, pero ese desplazamiento también puede degradar la regularidad superficial que hacía
que el SFT fuese tan estable. En una tarea parser-facing, esta regularidad no es secundaria: una pequeña caída de parseabilidad condiciona todas las métricas posteriores.

El experimento con \(\beta=0{,}30\) no modifica sustancialmente la lectura. Sobre las 42 muestras comparables, las salidas son prácticamente idénticas a las de \(\beta=0{,}10\) en 41
casos. Esto sugiere que, con el primer dataset de preferencias, aumentar \(\beta\) no proporciona una palanca clara de mejora. El siguiente movimiento metodológico no pasa tanto por
seguir elevando \(\beta\), sino por construir preferencias más discriminativas, especialmente con negativos duros que sean parseables y plausibles.

En el momento de redactar este capítulo, la segunda iteración DPO v2 se mantiene como experimento pendiente de cierre. Su configuración es deliberadamente más agresiva y parte del
checkpoint DPO v1, con un conjunto de preferencias más pequeño pero más estricto. Por esta razón, no debe anticiparse su conclusión. En la versión final de este capítulo, DPO v2 debe
añadirse solo después de completar su validación, comparándolo al menos con \texttt{serialized\_sft} y con DPO v1 sobre el mismo protocolo de métricas.

% TODO: incorporar aquí los resultados finales de DPO v2 cuando termine la validación.

%----------------------------------------------------------------------
\section{Evaluación final en test}
\label{sec:ch9_test_plan}
%----------------------------------------------------------------------

Hasta este punto, la mayor parte de la discusión se apoya en validación, porque ese split ha servido para desarrollar, comparar y diagnosticar las variantes experimentales. Para cerrar
el análisis comparativo sin contaminar el test, la evaluación final debe reservarse a un conjunto reducido de modelos ya seleccionados. La función del test no es elegir hiperparámetros
ni decidir entre variantes exploratorias, sino comprobar si la lectura obtenida en validación se mantiene en datos no usados durante el desarrollo.

La Tabla~\ref{tab:ch9_test_matrix} recoge la matriz propuesta para esa evaluación final. El baseline ya dispone de una ejecución en test y puede actuar como punto de partida. El SFT
serializado debe evaluarse porque es la referencia supervisada fuerte. La rama two-head debe representarse mediante el THM v1, que es el mejor resultado global de esa familia en
validación. Las variantes ordinales y posicionales quedan como análisis diagnóstico y como base para trabajo futuro, pero no como candidatas principales a test salvo que se quiera
confirmar explícitamente el fallo observado. DPO debe evaluarse solo cuando se haya decidido, usando validación, cuál es el checkpoint final que merece pasar a test.

\begin{table}[htbp]
\centering
\caption{Matriz de evaluación final en test. Las filas pendientes no deben usarse para selección de modelo, sino como confirmación de modelos previamente fijados en validación.}
\label{tab:ch9_test_matrix}
\resizebox{\textwidth}{!}{%
\begin{tabular}{lll}
\toprule
\textbf{Modelo} & \textbf{Estado} & \textbf{Función en la comparación} \\
\midrule
Baseline zero-shot & Ejecutado en test & Punto de partida sin ajuste supervisado \\
\texttt{serialized\_sft} & Pendiente de test & Referencia supervisada principal \\
THM v1 & Pendiente de test & Mejor representante validado de la rama con jerarquía explícita \\
DPO final & Pendiente de selección y test & Refinamiento post-SFT mediante preferencias automáticas \\
\bottomrule
\end{tabular}
}
\end{table}

Las métricas de esta evaluación deberían mantenerse alineadas con el resto del trabajo. En concreto, conviene agruparlas en cuatro bloques:

\begin{center}
\small
\begin{tabular}{ll}
\toprule
\textbf{Bloque} & \textbf{Métricas} \\
\midrule
Parseabilidad &
\texttt{yaml\_parse\_success\_rate} \\
&
\texttt{block\_parse\_success\_rate} \\
\addlinespace
Fidelidad estructural &
\texttt{average\_line\_text\_f1} \\
&
\texttt{average\_level\_exact\_match\_rate} \\
&
\texttt{average\_level\_mae} \\
\addlinespace
Adecuación al contenido &
\texttt{average\_prompt\_requirement\_f1} \\
&
\texttt{average\_semantic\_key\_f1} \\
\addlinespace
Dominio Kubernetes &
\texttt{average\_kubernetes\_domain\_validity\_score} \\
&
\texttt{kubernetes\_domain\_gate\_pass\_rate} \\
\bottomrule
\end{tabular}
\end{center}

Este conjunto permite comprobar no solo si el modelo genera YAML válido, sino también si conserva la jerarquía, responde al prompt y evita parte de los errores de dominio observados.

% TODO: sustituir esta sección por la tabla de resultados finales en test cuando estén disponibles las ejecuciones de SFT, THM v1 y DPO final.

%----------------------------------------------------------------------
\section{Discusión transversal}
\label{sec:ch9_cross_discussion}
%----------------------------------------------------------------------

Al poner juntas las tres líneas experimentales aparece una lectura más matizada que la que ofrecería cualquier métrica aislada. El SFT serializado es el sistema más fuerte porque
convierte una tarea jerárquica en una superficie textual muy regular y fácil de consumir por el parser. Esto no demuestra que el modelo haya aprendido una representación formal del
árbol YAML, pero sí demuestra que una serialización explícita de la jerarquía puede ser una estrategia práctica muy eficaz para generación estructurada.

La rama two-head, por el contrario, separa de manera más limpia la pregunta investigadora. Al sacar \texttt{level} de la superficie textual y predecirlo con una cabeza estructural,
permite estudiar si el modelo contiene información jerárquica aprovechable en sus estados latentes. El THM v1 muestra que esa separación puede producir un generador razonablemente
estable, aunque incapaz de recuperar niveles profundos por el desbalanceo de clases. Las variantes posteriores muestran el reverso del problema: cuando se fuerza al cabezal a entrar
en esa zona profunda mediante una formulación ordinal, la calibración de la trayectoria y la posición de la línea se vuelven críticas. La frontera entre generación de contenido,
predicción de nivel y reconstrucción parser-facing se convierte así en el punto frágil de la arquitectura.

DPO añade una tercera dimensión a la comparación. En lugar de cambiar la arquitectura, modifica la política aprendida a partir de preferencias automáticas. Su efecto inicial confirma
que las señales de dominio pueden mover el comportamiento del modelo, pero también que ese movimiento tiene coste. La mejora de algunas propiedades Kubernetes no compensa todavía las
pequeñas pérdidas de parseabilidad y fidelidad al prompt. En este sentido, DPO se interpreta mejor como una herramienta de refinamiento que como una sustitución directa del SFT.

La comparación global refuerza, por tanto, la necesidad de evaluar por niveles. Si solo se observara la parseabilidad YAML, el SFT serializado parecería cerrar el problema. Si solo se
observara la predicción de \texttt{level}, se perdería el hecho de que una pequeña desviación puede romper todo el documento. Si solo se observaran métricas de dominio, se ignoraría
que muchas buenas prácticas Kubernetes dependen de checks más estrictos que la sintaxis. El valor del análisis comparativo está precisamente en mantener juntas estas capas sin
confundirlas.

Este capítulo deja así preparada la lectura final del trabajo: los resultados no apuntan a una única técnica dominante en todos los sentidos, sino a una tensión metodológica entre
superficie serializada, señal estructural explícita y alineamiento automático. La interpretación de esa tensión, junto con las limitaciones y las líneas de continuidad, se desarrolla
en el capítulo siguiente.
